% problem-000293 2 蒙脱石纳米矿物提纯
\section*{2. 蒙脱石纳米矿物提纯}
\textit{下列为分问。}

\subsection*{2-1}
膨润⼟是⼀种天然纳⽶矿物材料，主要成分为蒙脱⽯。蒙脱⽯是层状结构的硅酸盐矿物，在⼀定地质条件下可形成厚度为纳⽶尺⼨的层状晶体。不考虑层间⽔和阳离⼦置换时，化学式为 $\mathrm { A l _ { 2 } S i _ { 4 } O _ { 1 0 } ( O H ) _ { 2 0 } }$ 将下⾯蒙脱⽯晶体结构图中与中⼼Al原⼦配位的O原⼦和OH原⼦团⽤短线连接起来（Al的配位多⾯体共边），指出配位多⾯体的空间构型。

（图：）

\subsection*{2-2}
⾮⾦属纳⽶矿物开采出来后，需经选矿、化学提纯才能应⽤。

\subsection*{2-2}
-1 当矿物中含有CuS、ZnS杂质时，将适量浓硫酸拌⼊矿粉，加热，充分搅拌，待反应完全后⽤⽔稀释、过滤、清洗，可将这些杂质分离除去。从环保和经济的⻆度，应尽量减少硫酸的⽤量，写出化学反应⽅程式。加⽔稀释的⽬的是什么？

\subsection*{2-2}
-2 当矿物中含有 ${ \mathrm { A s } } _ { 2 } { \mathrm { S } } _ { 3 } .$ $\mathrm { S n S _ { 2 } }$ 杂质时，加⼊硫化钠和少量氢氧化钠的混合溶液，加热搅拌，待反应完全后经过滤、消洗，可将杂质分离除去，写出主要化学反应⽅程式。为什么需加氢氧化钠溶液？

\subsection*{2-3}
若膨润⼟中含 $\mathrm { F e _ { 2 } O _ { 3 } }$ 时，⽩度就会受到影响，须进⾏漂⽩处理：在硫酸介质中⽤连⼆亚硫酸钠作漂⽩剂进⾏漂⽩；加⽔漂洗后加⼊氨基三⼄酸钠 $\left( \mathrm { N a } _ { 3 } \mathrm { A } \right)$ ，结合残留的亚铁离⼦，以维持产品的⽩度。

\subsection*{2-3}
-1 写出漂⽩过程的离⼦反应⽅程式。

\subsection*{2-3}
-2 写出氨基三⼄酸钠与亚铁离⼦结合⽣成的2:1型单核配合物的化学式及阴离⼦的结构式。
